\section{Each Element in Its Setting}
\label{sec:each-element}

Six of the ten elements are Scripture and four are the Church's own
composition. The chants come from five psalms: the Introit's verse, the
Gradual, the Alleluia, the Offertory and the Communion each take a verse or two
from a different psalm. The lessons are St Paul's letter to the Ephesians and
St Matthew's Gospel. The Introit's antiphon and the three orations are not
Scripture. The elements are set out below one by one, each in its own setting,
before the three readings of the whole.

\subsection{Introit: the Lord who speaks, and the psalm of Israel's history}

The antiphon puts its words in the Lord's mouth, \latin{dicit Dominus}, and
every clause of it is a promise: salvation for the people, a hearing for those
who cry \latin{de quacumque tribulatione}, from any tribulation whatever, and
lordship \latin{in perpetuum}. The words stand
in no single verse of Scripture. The continuation of Dom Prosper Guéranger's
\work{The Liturgical Year} by Dom Lucien Fromage, in its English of 1883, says
of the antiphon that
it ``is made up of several passages of Holy Writ, without being exactly that of
any one of them''. The psalter's promises lie close to it, God's ``I am thy
salvation'' to the soul (Ps 34:3) and ``He shall cry to me, and I will hear
him: I am with him in tribulation'' (Ps 90:15), and the Offertory's own psalm
has the same prayer in the psalmist's voice, ``In what day soever I shall call
upon thee, hear me'' (Ps 137:3). The antiphon's word \latin{tribulatio} comes
back once, in the Offertory, and nowhere else in the formulary.

The verse is the first of Ps~77, titled in the Vulgate \latin{Intellectus
Asaph}. It is one of the longest psalms of the psalter, and it tells Israel its
own history: the fathers have told their children what the Lord did, ``That
another generation might know them'' and ``may not become like their fathers, a
perverse and exasperating generation'' (vv.~3--8). The psalm then recounts the
wonders in Egypt, the dividing of the sea, the wilderness and the entry into
the land, and ends with the Lord's choice of Juda, of Mount Sion and of his
servant David, taken ``from the flocks of sheep \dots\ To feed Jacob his
servant'' (vv.~12--13, 52--54, 67--72). The verse that follows the one the
Introit sings is ``I will open my mouth in parables'' (v.~2), which St Matthew
cites as fulfilled in Jesus's speaking in parables (Mt 13:35). The Gospel of
this Mass opens with Jesus speaking \latin{in parabolis}.

St Augustine asks of the psalm's first words, ``Whom may we suppose to be here
speaking, but God?'' (\work{Enarrationes in Psalmos} 77, 2). The antiphon has
already named the speaker, \latin{dicit Dominus}. The medieval liturgical
commentators Honorius Augustodunensis (\work{Gemma animae} IV.86) and Sicard of
Cremona (\work{Mitrale} VIII.19) set the antiphon in the return from the
Babylonian exile: after Cyrus let the captive people go, Zorobabel and the
priests proclaimed to them the Lord's word that he was their salvation.
Bl.\ Ildefonso Schuster, the Benedictine liturgist and Cardinal Archbishop of
Milan, reads it as the Lord's own assurance against every human helper: ``It
is useless, therefore, to put one's trust in the arm of man which must perish
and turn to dust, whereas no one has ever called upon the Lord for help and
been disappointed'' (\work{The Sacramentary} III, p.~171).

\subsection{Collect: freedom for the things that are God's}

The Collect is built on a single contrast between what hinders and what is set
free. It asks God, \latin{propitiatus}, to shut out \latin{universa nobis
adversantia}, all that is adverse, and it gives the purpose in a clause that
names freedom twice over, \latin{mente et corpore pariter expediti} and
\latin{liberis mentibus}. The participle \latin{expediti} means disencumbered,
set loose for action, and it is said of mind and body together; the second
phrase means with free minds. What the freed are to carry out is named only as
\latin{quae tua sunt}, the things that are God's. The freedom asked for is a
freedom for action, and its object is God's.

Schuster reads the prayer as a petition for the order God set between body and
soul: ``The body, when healthy and strong, must act obediently under the rule
of the soul, which on its part is raised up on high by grace, to God'', so
that ``the whole man is directed gently and without difficulty or hindrance
towards God as his final end''. ``Glorious and beautiful indeed is this liberty
of the sons of God, for it is the result of order, harmony, and the due
subjection of the creature to the Creator'' (pp.~171--172). Honorius, reading
the same Collect of the people released from captivity, has them ask that
everything adverse to them be shut out.

Two words of the Collect return in the Postcommunion. \latin{Expediti} comes
back as \latin{expediat}, and the things that are God's, \latin{quae tua sunt},
come back as \latin{tuis \dots\ mandatis}, God's commandments, to which the Mass
at its end asks to cleave.

\subsection{Epistle: the new man, and four acts}

The lesson begins in the middle of a sentence. Paul has just told the
Ephesians ``To put off, according to former conversation, the old man, who is
corrupted according to the desire of error'' (Eph 4:22), and the verses the
Missal reads complete it: be renewed in the spirit of your mind, and put on
the new man, created according to God ``in justice and holiness of truth''
(4:23--24). The chapter from which the lesson is taken begins with Paul, ``a
prisoner in the Lord'', beseeching his readers ``that you walk worthy of the
vocation in which you are called'' (4:1), \latin{ut digne ambuletis vocatione,
qua vocati estis}.

After the general command come four particular ones, and each names an act of
the old man to be put off and an act of the new to be put on. Lying gives way
to truth spoken to the neighbour, with a reason: ``For we are members one of
another'' (4:25). Anger is allowed its first motion and refused its duration:
``Be angry: and sin not. Let not the sun go down upon your anger'' (4:26). The
devil is to be given no place (4:27). And the thief is to become a worker who
gives: ``let him labour, working with his hands the thing which is good, that
he may have something to give to him that suffereth need'' (4:28). The verses
after the lesson forbid evil speech, warn ``grieve not the holy Spirit of God:
whereby you are sealed unto the day of redemption'', and end with forgiveness,
``even as God hath forgiven you in Christ'' (4:29--32).

The chief commentators on the Epistle did not all read its words as the Missal
prints them. St Jerome's text of verse 23 had \latin{spiritu sensus vestri}
where the Missal has \latin{mentis}, and his verse 28 ended \latin{ut habeat
communicare ei qui indiget}, that he may have something to share with one who
is in want. St Thomas Aquinas's printed text divides the verses differently
from the Clementine, so that his twenty-seventh verse is the Missal's
twenty-eighth. Neither difference touches the sense of the lesson.

Schuster reads the lesson in one sentence: ``St Paul sets before us Christian
sanctity under the symbol of a new garment, which we are to put on. This
garment is Jesus Christ himself, with his divine qualities and feelings''; the
faithful, ``being members of one and the same mystical body'', owe one another
``charity, sincerity, trustfulness, and compassion'' (p.~172). Honorius reads the
new man as the condition of the exiles' return: they are warned to put on the
new man through justice so that, God leading them, they may go back to their
homeland. Sicard, too, says that the new man must be put on.

\subsection{Gradual: prayer as incense, hands lifted at evening}

Ps~140 is titled in the Vulgate \latin{Psalmus David}, and it is an evening
prayer. It begins with a cry, ``I have cried to thee, O Lord, hear me'' (v.~1),
and the Gradual's verse compares the prayer to the incense and the evening
sacrifice of Israel's worship: ``Let my prayer be directed as incense in thy
sight; the lifting up of my hands, as evening sacrifice'' (v.~2). The next
verses ask for a guard on the mouth and a heart kept from ``evil words; to make
excuses in sins'' (vv.~3--4). The Missal adds \latin{Domine} to the verse, and
its verb \latin{dirigatur}, may it be directed or set straight, comes back in
the Communion's \latin{dirigantur}.

The same verse returns in the Order of Mass. When the Mass is celebrated with
incense, the priest incenses the altar at the Offertory saying Ps 140:2--4,
\latin{Dirigatur, Domine, oratio mea, sicut incensum, in conspectu tuo:
elevatio manuum mearum sacrificium vespertinum} (1962 \work{Ordo Missae},
no.~1037). At a low Mass there is no incense and the verse is not said a second
time.

Honorius has the people released from exile, already on the road, invoke the
Lord with this verse. Schuster reads it of every Christian offering in this
life, which is always an evening sacrifice for him because it is made in the
twilight of faith (p.~172). St Augustine reads the verse's evening sacrifice of
the Lord's Passion, and St Hilary of the works of mercy.

\subsection{Alleluia: praise, invocation and proclamation}

Ps~104, St Augustine observes, ``is the first of those to which is prefixed
the word Allelujah'' (\work{Enarrationes} 104, 1). Its first verse holds three
commands in a row: give glory, call upon his name,
declare his deeds among the Gentiles. The psalm then recounts the covenant with
Abraham, Isaac and Jacob, the story of Joseph, the plagues of Egypt, the
Exodus and the gift of the land, and it gives the reason for all of it in its
last verse: ``That they might observe his justifications, and seek after his
law'' (v.~45). The Communion of this Mass asks that the ways of the one who prays
be directed \latin{ad custodiendas iustificationes tuas}, to keep those same
justifications. The psalm's first verses stand also in the First Book of
Paralipomenon, in the praise that David appointed Asaph and his brethren to
sing on the day the ark was set in the tent he had pitched for it (1 Par 16:1,
7--8).

Augustine also noticed that other Latin copies had another verb for the last
clause: \latin{Annuntiate inter gentes opera ejus; vel potius, ut de graeco ad
verbum exprimatur, quod et alii latini codices habent, Evangelizate in gentibus
opera ejus}, in the English of the Nicene and Post-Nicene Fathers, ``or rather,
to translate literally from the Greek, as other Latin copies too have it,
`Preach the Gospel of His works among the Gentiles'\,'' (\work{Enarrationes}
104, 1). He heard in the verse a command given in prophecy to the evangelists. St Robert Bellarmine reads the verse as the voice of a lover:
``The true lover does not wish the praise and knowledge of his beloved should
be confined to himself, but wishes that many, nay even all, should know her
perfections and praise them'', so that the Gentiles, too, ``from a knowledge of
his works, may begin to know, praise, and invoke their Creator''
(\work{Commentary on the Psalms}, on Ps 104:1). Schuster names those who carry
out the command: ``The apostles, and after them the bishops and pastors of
souls, have always looked upon this duty of preaching the Gospel as the most
important of all'' (p.~172).

\subsection{Gospel: the king's marriage feast}

The Missal's opening names the parable's hearers, \latin{principibus sacerdotum
et pharisaeis}, and so sets it where St Matthew sets it, in the temple at
Jerusalem after the entry into the city. There the chief priests and ancients
of the people have asked Jesus by what authority he acts (Mt 21:23). He has
answered with two parables, the two sons and the wicked husbandmen, has
quoted the stone the builders rejected, and has said, ``the kingdom of God
shall be taken from you and shall be given to a nation yielding the fruits
thereof'' (21:43). ``And when the chief priests and Pharisees had heard his
parables, they knew that he spoke of them'' (21:45). The parable of the
marriage feast is his third answer, and after it the Pharisees go out to take
counsel ``how to insnare him in his speech'' (22:15).

The parable moves in two halves. In the first (vv.~2--10) a king sends his
servants twice to those invited to his son's wedding. The first time ``they
would not come''. The second time the servants announce that the dinner is
prepared, ``my beeves and fatlings are killed, and all things are ready'', and
the invited make light of it and go off, ``one to his farm and another to his
merchandise'', while the rest seize the servants, treat them shamefully and
kill them. The king is angry, sends his armies, destroys the murderers and
burns their city. Then he declares the invited unworthy and sends his servants
\latin{ad exitus viarum}, to the partings or ends of the roads, to call
everyone they find; they gather ``both bad and good'', and the hall is filled.
In the second half (vv.~11--14) the king comes in to see those at table and
finds one man not clothed in a wedding garment. He calls him \latin{Amice},
friend, asks how he came in, and receives no answer, \latin{At ille
obmutuit}. The man is bound hand and foot and cast into the exterior darkness,
and the parable ends with a saying that reaches beyond the story: ``For many
are called, but few are chosen.''

St Luke has a parable of a great supper whose invited guests all excuse
themselves and whose host fills the house from the streets and lanes, the
highways and hedges (Lk 14:16--24). St Gregory the Great opens his homily on
Matthew's parable by asking whether the two are one lesson, and answers that
Luke's supper signifies the eternal and last banquet and Matthew's wedding the
present Church, because at Matthew's wedding one guest who came in is put out
again, and at the eternal banquet no one who has once come in goes out
(\work{Homiliae in Evangelia} 38, 1). He preached that homily to the people in
the basilica of the martyr St Clement at Rome, as its heading records;
Schuster adds that the occasion is not known (p.~172).

The Missal's \latin{Tunc dixit rex}, where the Clementine Vulgate has
\latin{dicit}, is the reading of St Gregory's and St Jerome's texts of the
verse. The continuation of \work{The Liturgical Year} reports that this Gospel
``has given to the present Sunday the name of the Sunday of the invited to the
marriage''.

\subsection{Offertory: walking in tribulation}

Ps~137, titled \latin{Ipsi David}, is a thanksgiving. The psalmist praises God
with his whole heart, ``for thou hast heard the words of my mouth'', sings
``in the sight of the angels'' and worships ``towards thy holy temple''
(vv.~1--2), asks to be heard on whatever day he calls (v.~3), and prays that
``all the kings of the earth give glory to thee'' (v.~4). The Offertory's verse turns from past deliverance to trust in what is
to come. Where the Vulgate says that God has stretched forth his hand and that
his right hand has saved, the Missal's chant sets both verbs in the future, and
the future already stood in the Pustet Missal of 1862. St Augustine and St
Hilary read the verse with the Vulgate's past, and St John Chrysostom's Greek
text of the verse has the past as well, though Chrysostom reports two other
Greek versions that read the future. The chant accompanies the offering of
the gifts.

Schuster finds in it the hour at which God acts: ``As long as man continues to
act, the Lord would seem to withhold his assistance, but when all human hope is
vain, then it is that God most often begins his work of salvation'', so that
``the hour in which God hears our prayer is the darkest and saddest of our
trial'' (p.~173). Honorius has the people returned from exile give thanks in
this chant for their return.

\subsection{Secret: gifts before the eyes of God's majesty}

The Secret is short and asks one thing, that the gifts offered \latin{oculis
tuae maiestatis} be \latin{salutaria nobis}, saving to us. The image of offering
before God's eyes runs through this part of the Mass. The Gradual prayed that
the psalmist's prayer rise \latin{in conspectu tuo}, in God's sight. At the
offering of the chalice the priest asks that it rise \latin{in conspectu
divinae maiestatis tuae} with an odour of sweetness (\work{Ordo Missae},
no.~1032). And at the end of every Mass, low or solemn, in the prayer
\latin{Placeat}, he asks that the sacrifice \latin{quod oculis tuae maiestatis
indignus obtuli}, which he unworthy has offered before the eyes of God's
majesty, may be acceptable to God and \latin{propitiabile} for himself and all
for whom he offered it (no.~1137). The Secret and the \latin{Placeat} share the three words \latin{oculis tuae
maiestatis}.

Schuster reads the Secret of the Eucharist and of those who share in it, as a
prayer that the oblation be a pledge of salvation to them and that they bring
it the dispositions its effect in each of them requires (p.~173).

\subsection{Communion: a command and a wish}

The Communion's two verses open the long psalm of the law, Ps~118, in its first
strophe, which the Vulgate heads with the Hebrew letter \latin{Aleph}. The
psalm begins with a beatitude, ``Blessed are the undefiled in the way, who walk
in the law of the Lord'' (v.~1). Its fourth verse states a command and its
fifth answers with a wish: \latin{Tu mandasti mandata tua custodiri nimis:
utinam dirigantur viae meae, ad custodiendas iustificationes tuas}. The word
\latin{nimis} does not mean ``too much'', as St Augustine explains, but ``very
much'' (\work{Enarrationes} 118, Aleph 4). St Robert Bellarmine hears the
first verse as a statement of the lawgiver's authority: God has commanded
``not by way of advice, but by strict precept'' (on Ps 118:4). The
second verse turns the command into a prayer. Its \latin{dirigantur} answers
the Gradual's \latin{dirigatur}: the evening prayer asked to be set straight
like incense, and the Communion asks that the ways of the one who prays be set
straight towards what God commands.

Schuster hears joy in the wish: ``by the exclamation \latin{utinam}, he
signifies his joy and delight in the service of God'' (p.~173). Honorius has
the returning people profess in the Communion that they keep the law of God,
and ask for the power to fulfil it.

\subsection{Postcommunion: the healing work}

The Postcommunion names God's work by one phrase,
\latin{tua \dots\ medicinalis operatio}, thy healing working, and asks two
things of it. The first is the Collect's freedom again, now as freedom from
something inside us rather than from adversities without: \latin{a nostris
perversitatibus clementer expediat}, in the 1861 English ``mercifully free us
from our perverseness''. The second is the Communion's keeping again:
\latin{tuis semper faciat inhaerere mandatis}, where the 1861 English has
``make us always obedient to thy commandments'' and the Latin's verb,
\latin{inhaerere}, means to cleave or cling.
The Mass that began with God promising to be his people's Lord \latin{in
perpetuum} ends with a prayer to hold to his commandments \latin{semper}.

Schuster reads the prayer of the Eucharist, whose healing grace it asks to
free us from our perversity, and he names the sickness that needs the
medicine: ``our infirmity is caused by the poisoning of our nature through the
fatal apple of Eden'' (p.~174).
